%===========================================================================
% APPENDIX: Identification of PROGRESA Transfers in the ENIGH Survey
% For: "The Effect of PROGRESA on Adult Mortality"
% Journal: Journal of Development Economics
%===========================================================================

\clearpage
\appendix

\section{Constructing PROGRESA Transfer Variables from the ENIGH}
\label{app:enigh_progresa}

This appendix documents the data construction of two distinct variables used to
capture PROGRESA exposure in the Encuesta Nacional de Ingresos y Gastos de los
Hogares (ENIGH): (i) the \textit{income amount} of PROGRESA transfers received
by a household or individual (\texttt{progresa\_hh}, \texttt{progresa\_ind}),
and (ii) a binary \textit{beneficiary flag} indicating whether any household
member holds a PROGRESA scholarship (\texttt{progresa\_benef\_hh},
\texttt{progresa\_benef\_ind}). Both variables are used as \textit{outcome
variables} in the mechanisms analysis (Section~\ref{sec:mechanisms}); they are
\textit{not} used to assign program treatment, which is drawn from
administrative enrollment records (see \texttt{aamr\_011326.do} and
Section~\ref{sec:data}).

The central data challenge is that PROGRESA did not have a dedicated income
code in the ENIGH until 2002. The program launched in the second half of 1997
and began disbursing transfers in Spring 1998, but INEGI had not yet adapted
the survey instrument to capture it as a named program. We use five ENIGH waves
covering the pre- and immediate post-program period (1992, 1994, 1996, 1998,
2000), plus three post-2002 waves (2002, 2004/2005, and 2006 for robustness).
The construction strategy differs materially across periods.

%--------------------------------------------------------------------------
\subsection{ENIGH Survey Instrument Evolution for Program Transfers}
\label{app:enigh_evolution}
%--------------------------------------------------------------------------

Table~\ref{tab:enigh_pcodes} summarises how PROGRESA/Oportunidades transfers
are recorded in the ENIGH income questionnaire across waves, and where other
government institutional transfers reside in the same questionnaire.

\begin{table}[h!]
\centering
\caption{ENIGH Income P-Codes for Transfers, by Wave}
\label{tab:enigh_pcodes}
\small
\begin{tabular}{lllll}
\hline\hline
Wave & PROGRESA code & Government transfers & Non-gov.\ transfers & Notes \\
\hline
1992 & --- & P025 & P026 & Program did not exist \\
1994 & --- & P026 & P027 & Program did not exist \\
1996 & --- & P026 & P027 & Program did not exist \\
1998 & P031 (implicit) & P031 & P032 & SEDESOL annotation in 2000 \\
     &                  &      &      & manual; P031 = becas/donativos \\
     &                  &      &      & de instituciones (all types) \\
2000 & P031 (explicit) & P031 & P032 & Enumerator manual instructs: \\
     &                  &      &      & ``PROGRESA se registra en este \\
     &                  &      &      & rengl\'on'' \\
2002 & P046 (own code) & P043 & P042 & First wave with dedicated \\
     &                  &      &      & PROGRESA/Oportunidades code \\
2004/2005 & P059 (own code) & P056 & P055 & Fully disaggregated \\
2006 & P059 (own code) & P056 & P055 & Same as 2004/2005 \\
\hline\hline
\end{tabular}
\end{table}

\noindent
The 2000 ENIGH enumerator manual (\textit{Manual del Entrevistador}, INEGI
2000) contains an explicit SEDESOL annotation embedded in the instructions for
P031:

\begin{quote}
\textit{``La Secretar\'ia de Desarrollo Social (SEDESOL) ha implementado un
mecanismo de ayuda para las familias de escasos recursos (Programa Progresa).
Cuando los ingresos monetarios sean otorgados por este programa el cual es un
apoyo econ\'omico que se proporciona en efectivo en la modalidad de beca o
donativo, se registran en este rengl\'on al miembro del hogar responsable del
menor.''}
\end{quote}

This annotation confirms that PROGRESA cash disbursements---covering the
educational, nutritional, and health attendance components paid to the
\textit{vocal} (typically the mother)---were directed to P031 (\textit{Becas y
Donativos provenientes de Instituciones}) for the 2000 wave. The same P-code
numbering applied in 1998, although the 1998 enumerator manual did not yet
contain the SEDESOL annotation. A parallel variable in the ENIGH population
file (\texttt{pob.dta}) captures scholarship type at the individual level:
\texttt{PROP\_BECA} in 2000, with code 5 labelled \textit{Becas de PROGRESA}
in the survey catalog. This individual-level flag does not exist in 1998 or
earlier waves.

%--------------------------------------------------------------------------
\subsection{Construction of the PROGRESA Transfer Amount (Income Variable)}
\label{app:enigh_amount}
%--------------------------------------------------------------------------

\subsubsection{Waves 1992, 1994, 1996}

PROGRESA did not exist prior to 1997. For these pre-program waves we set
\texttt{progresa\_d = 0}, so \texttt{progresa\_ind} and \texttt{progresa\_hh}
are zero for all observations. These waves serve as the clean pre-treatment
baseline in the mechanisms regressions.

\subsubsection{Wave 2002, 2004/2005, 2006}

Starting in 2002, PROGRESA/Oportunidades has a dedicated P-code: P046 in 2002
and P059 in 2004--2006. The income variable is constructed directly from these
codes with no imputation. We define the income indicator \texttt{progresa\_d}
as:

\begin{align*}
\texttt{progresa\_d}_{i} &=
  \mathbf{1}\bigl[\text{code} = \texttt{P046}\bigr] \quad \text{(2002)} \\
\texttt{progresa\_d}_{i} &=
  \mathbf{1}\bigl[\text{code} = \texttt{P059}\bigr] \quad \text{(2004, 2005, 2006)}
\end{align*}

and the household-level variable is the sum of all individual P-code amounts
within a household (reported quarterly, converted to monthly: \texttt{ing\_tri
/ 3}).

\subsubsection{Waves 1998 and 2000: Imputation from the 2002 Share}
\label{app:enigh_imputation}

For 1998 and 2000, P031 collapses PROGRESA transfers together with all other
government institutional scholarships and grants (CONAFE stipends, municipal
scholarship programmes, etc.). There is no way to disentangle these ex ante at
the respondent level. We use the following imputation, applied after appending
all waves:

\begin{equation}
\widehat{\texttt{progresa}}_{i,t} =
  \underbrace{\left(\texttt{benef\_don\_gob}_{i,t}
    + \texttt{benef\_don\_non\_gob}_{i,t}\right)}_{\text{P031} + \text{P032
    amounts}}
  \times \hat{\omega}_{2002}
\label{eq:progresa_share}
\end{equation}

\noindent for $t \in \{1998, 2000\}$, where

\begin{equation}
\hat{\omega}_{2002} =
  \frac{\overline{\texttt{progresa}}_{2002}}
       {\overline{\texttt{progresa}}_{2002}
        + \overline{\texttt{benef\_don\_gob}}_{2002}
        + \overline{\texttt{benef\_don\_non\_gob}}_{2002}}
\label{eq:share_2002}
\end{equation}

is PROGRESA's share of the combined P046 + P043 + P042 pool in 2002, computed
as the ratio of sample means at the national level. \texttt{benef\_don\_gob}
corresponds to government institutional scholarships/grants (P031 in 1998/2000;
P043 in 2002) and \texttt{benef\_don\_non\_gob} corresponds to
non-governmental institutional scholarships/grants (P032 in 1998/2000; P042 in
2002).

The logic is that the composition of the pooled transfer category in 2002---the
first year PROGRESA appears with its own code and therefore the first year in
which the shares can be identified---proxies for its composition in the earlier
waves. This is valid under the assumption that PROGRESA's share of total
institutional transfers within the target population did not change drastically
between 1998/2000 and 2002. While program enrollment expanded between 1998 and
2002, the composition of \textit{other} institutional transfers (CONAFE,
municipal programmes) was also relatively stable, supporting this assumption.
The imputation applies to both individual- and household-level variables
identically.

The Stata implementation is:

\begin{verbatim}
* Compute 2002 share from sample means
sum benef_don_non_gob_ind if year == 2002
local m_benef_non_gob = `r(mean)'
sum benef_don_gob_ind   if year == 2002
local m_benef_gob = `r(mean)'
sum progresa_ind        if year == 2002
local m_progresa = `r(mean)'

local share_2002 = `m_progresa' / ///
    (`m_progresa' + `m_benef_gob' + `m_benef_non_gob')

* Impute 1998 and 2000 PROGRESA amounts
replace progresa_ind = ///
    (benef_don_non_gob_ind + benef_don_gob_ind) * `share_2002' ///
    if inlist(year, 1998, 2000)
\end{verbatim}

%--------------------------------------------------------------------------
\subsection{Alternative Proxies and Robustness}
\label{app:enigh_alternatives}
%--------------------------------------------------------------------------

The baseline imputation uses a single national share and pools government (P031)
and non-governmental (P032) transfers in the denominator. Two alternative
proxies address these limitations.

\paragraph{Alternative 1 (Alt1): State-level share, same pool.}
The ENIGH is designed to be representative at the state level. Conditioning the
share on state rather than national averages captures geographic variation in
PROGRESA rollout. For state $e$:

\begin{equation}
\hat{\omega}_{e,2002}^{\text{Alt1}} =
  \frac{\overline{\texttt{progresa}}_{e,2002}}
       {\overline{\texttt{progresa}}_{e,2002}
        + \overline{\texttt{benef\_don\_gob}}_{e,2002}
        + \overline{\texttt{benef\_don\_non\_gob}}_{e,2002}}
\end{equation}

where the overlines denote state-specific means. States with no positive 2002
observations fall back to the national mean $\hat{\omega}_{2002}$.

\paragraph{Alternative 2 (Alt2): State-level share, government-transfers
pool only.}
The SEDESOL annotation instructs enumerators to record PROGRESA under the
government channel (P031), not the non-governmental channel (P032). The
non-governmental code (P032) captures private gifts, remittances from
institutions, and civil-society transfers---flows unrelated to PROGRESA. Alt2
restricts the denominator to the government channel and scales only the
government-source amount in 1998/2000:

\begin{equation}
\hat{\omega}_{e,2002}^{\text{Alt2}} =
  \frac{\overline{\texttt{progresa}}_{e,2002}}
       {\overline{\texttt{progresa}}_{e,2002}
        + \overline{\texttt{benef\_don\_gob}}_{e,2002}}
\label{eq:alt2}
\end{equation}

\begin{equation}
\widehat{\texttt{progresa}}_{i,t}^{\text{Alt2}} =
  \texttt{benef\_don\_gob}_{i,t}
  \times \hat{\omega}_{e,2002}^{\text{Alt2}}
  \quad t \in \{1998, 2000\}
\end{equation}

Alt2 is theoretically cleaner because it matches the channel documented in the
enumerator manual. However, it drops non-governmental institutional transfer
amounts from the 1998/2000 imputed total, which may understate amounts in
localities where PROGRESA was disbursed through channels recorded under P032
(e.g., via NGO intermediaries in some early rollout areas).

%--------------------------------------------------------------------------
\subsection{Construction of the Beneficiary Flag}
\label{app:enigh_benef_flag}
%--------------------------------------------------------------------------

The binary beneficiary variable (\texttt{progresa\_benef\_ind}) captures
whether a household member is identified as a current scholarship recipient in
the ENIGH population file (\texttt{pob.dta}). This variable exists in the
population module and is distinct from the income-module transfer amounts. It
is constructed as follows by wave:

\begin{table}[h!]
\centering
\caption{Construction of \texttt{progresa\_benef\_ind} by ENIGH Wave}
\label{tab:benef_flag}
\small
\begin{tabular}{lll}
\hline\hline
Wave & Population variable & Condition \\
\hline
1992 & --- & = 0 (program did not exist) \\
1994 & --- & = 0 (program did not exist) \\
1996 & --- & = 0 (program did not exist) \\
1998 & --- & = 0 (no PROGRESA entry in \texttt{pob.dta}) \\
2000 & \texttt{PROP\_BECA} & = 1 if \texttt{PROP\_BECA} == 5 \\
     &                     & (catalog: \textit{Becas de PROGRESA}) \\
2002 & \texttt{BECA}       & = 1 if \texttt{BECA} == 5 \\
2004/2005 & \texttt{BECAS} & = 1 if \texttt{BECAS} == 1 \\
2006 & \texttt{BECA}       & = 1 if \texttt{BECA} == 1 \\
\hline\hline
\end{tabular}
\end{table}

\noindent
The household-level flag (\texttt{progresa\_benef\_hh}) equals one if any
household member satisfies the individual condition:

\begin{verbatim}
bys FOLIO: egen progresa_benef_hh = max(progresa_benef_ind)
\end{verbatim}

\paragraph{Coverage limitation for 1998.}
Although PROGRESA began disbursing transfers in Spring 1998, the 1998 ENIGH
population file does not contain a scholarship-type variable that identifies
PROGRESA recipients. We therefore set \texttt{progresa\_benef\_ind = 0} for
all 1998 respondents. This creates a mechanical zero in 1998 for the beneficiary
flag even among true PROGRESA recipients. Given that we use this variable as an
outcome variable rather than a treatment indicator, the practical consequence is
that regressions using the beneficiary flag as a dependent variable are
interpreted as showing take-up visibility \textit{in the survey}, and 1998
should be interpreted cautiously.

\paragraph{Coverage limitation for 2000.}
\texttt{PROP\_BECA == 5} identifies individuals who report a \textit{Beca de
PROGRESA} scholarship in the education module. PROGRESA's educational component
was disbursed to the mother/household head on behalf of the enrolled child;
the child is the scholarship holder. The health and nutrition transfer
components---disbursed unconditionally conditional on clinic attendance---are
not captured by \texttt{PROP\_BECA}. Thus, the 2000 beneficiary flag
understates true beneficiary status: it records households with at least one
school-age child receiving the educational scholarship, but misses households
that received only nutrition/health transfers (e.g., with children under 3 or
pregnant/lactating women but no school-age children). The beneficiary flag is
therefore a \textit{lower bound} on true beneficiary status in 2000.

%--------------------------------------------------------------------------
\subsection{Identifying Assumptions and Known Limitations}
\label{app:enigh_assumptions}
%--------------------------------------------------------------------------

The validity of the income-amount variable (\texttt{progresa\_ind}) for 1998
and 2000 rests on three key assumptions, each with associated caveats.

\paragraph{Assumption 1: Stable composition of institutional transfers.}
We assume $\hat\omega_{2002}$ (PROGRESA's fraction of the pooled institutional
transfer category) is a valid proxy for $\omega_{1998}$ and $\omega_{2000}$.
This holds if other components of P031 (government side) did not grow
disproportionately faster than PROGRESA between 1998 and 2002. Given that
PROGRESA was by far the largest targeted government transfer programme in Mexico
in this period, the assumption is plausible in highly marginalized rural
municipalities (our sample restriction), where PROGRESA dominated the
institutional transfer category.

\paragraph{Assumption 2: No attenuation from over-reporting of 2002 PROGRESA.}
\citet{hernandez2007} document that Oportunidades income reported in the 2002
ENIGH grew by more than 100 percent relative to 2000, while administrative
budget data show only 59 percent growth. The authors attribute this partly to
an oversample of Oportunidades households funded by SEDESOL for the 2002
round, which inflated both the numerator and the denominator in
Equation~\eqref{eq:share_2002}. Crucially, the effect on $\hat\omega_{2002}$
is ambiguous in sign: if PROGRESA income is over-reported more than other
government transfers in 2002, $\hat\omega_{2002}$ is inflated, causing upward
bias in the imputed 1998/2000 amounts. We cannot directly correct for this
without knowing the true over-report factor for PROGRESA versus other transfers;
results using \texttt{progresa\_ind} as an outcome should be read as reflecting
the transfer proxy with potential upward bias in levels for 1998 and 2000. Our
preferred robustness check (Alt2) reduces but does not eliminate this concern.

\paragraph{Assumption 3: No P031/P032 contamination in the analytical sample.}
Even after imputation, \texttt{progresa\_ind} in 1998/2000 may include
non-PROGRESA institutional transfers (CONAFE tutor stipends, municipal
scholarship programmes). In our sample of highly marginalized municipalities
(marginalization grades 4--5 in 1990), PROGRESA was the only large-scale
government transfer programme targeting education and nutrition, so
contamination is likely minimal relative to the national sample. We cannot
verify this directly.

%--------------------------------------------------------------------------
\subsection{Role in the Empirical Analysis}
\label{app:enigh_role}
%--------------------------------------------------------------------------

Both variables serve exclusively as \textit{dependent variables} in the
mechanisms analysis. We do not use ENIGH-based PROGRESA measures to assign
treatment. Treatment is assigned based on administrative records of
PROGRESA/Oportunidades enrollment at the municipality level, which provide
clean beneficiary counts for each year 1997--2018 (see
Section~\ref{sec:data_prog}).

The income-amount variable (\texttt{progresa\_hh}) appears in Table~X
(mechanisms panel: transfers received) and serves to verify that PROGRESA
transfers are visible in the survey among households in high-intensity
municipalities---a standard plausibility check for mechanisms analysis using
survey data. The beneficiary flag (\texttt{progresa\_benef\_hh}) is used as an
alternative LHS variable to assess extensive-margin take-up (whether a
household has any beneficiary), complementing the intensive-margin amount.

The use of PROGRESA variables as outcomes rather than treatment indicators
means that imprecision or attenuation in the proxy affects our ability to
\textit{detect} mechanisms through transfers---it creates measurement error in
the dependent variable, which inflates standard errors but does not bias
coefficients. Concerns about over-reporting in 2002 would affect levels
comparisons between years but not the within-municipality difference-in-
differences estimator which absorbs municipality fixed effects.

%--------------------------------------------------------------------------
\subsection{Cross-Wave Comparability and Deflation}
\label{app:enigh_deflation}
%--------------------------------------------------------------------------

All income amounts are reported quarterly (\texttt{ing\_tri}) and converted to
monthly by dividing by three (\texttt{ing\_tri / 3}). Amounts are then
converted to real 2025 USD using wave-specific CPI deflators and the
contemporaneous MXN/USD exchange rate (variable \texttt{inf\_ex\_rate} in the
code), applied uniformly across all income variables. The 1992 wave reports
amounts in old pesos (prior to the 1993 redenomination at 1000:1) and is
rescaled accordingly (\texttt{GAS\_TRI / 1000} before all other operations).

%--------------------------------------------------------------------------
\subsection{Literature Context}
\label{app:enigh_lit}
%--------------------------------------------------------------------------

The challenge of identifying PROGRESA beneficiaries in pre-2002 ENIGH data is
widely acknowledged. \citet{hernandez2007} provide the definitive documentation
of cross-wave comparability issues, noting that the number of income questions
grew from 36 in 2000 to 48 in 2002 to 61 in 2004 and that SEDESOL partially
funded the 2002 ENIGH specifically to oversample Oportunidades households.
Researchers requiring clean PROGRESA identification therefore typically adopt
one of three strategies: (i) use the PROGRESA experimental evaluation surveys
(ENCEL/ENCASEH), which provide clean beneficiary status from administrative
records \citep{gertler2004,parker2005,skoufias2001}; (ii) merge administrative
enrollment records at the municipality level directly, bypassing the survey
altogether \citep{barham2013}; or (iii) restrict ENIGH-based mechanisms
analysis to 2002 and later waves where PROGRESA has a dedicated code
\citep{parker2008}. Our approach---using 1998/2000 with the 2002-share
imputation as a \textit{dependent} variable check---is most similar in spirit
to strategy (ii), with the survey serving as a validity check rather than as
the primary source of treatment variation.

\citet{barham2013} is the closest existing study to our design: they examine
the impact of PROGRESA on elderly mortality using a municipality-level
difference-in-differences strategy with administrative enrollment data and vital
statistics, without relying on the ENIGH for treatment assignment. Our
extension uses the ENIGH mechanisms analysis to shed light on the economic
channels (income, food consumption, health expenditure) through which the
observed mortality effects operate.

%--------------------------------------------------------------------------
% REFERENCES (include in main bibliography)
%--------------------------------------------------------------------------
% @article{hernandez2007,
%   author  = {Hern\'andez Meza, Gonzalo and others},
%   title   = {Los problemas de comparabilidad de las {ENIGH} y su efecto
%              en la medici\'on de la pobreza},
%   journal = {Papeles de Poblaci\'on},
%   year    = {2007},
%   number  = {51},
%   pages   = {95--125}
% }
%
% @article{gertler2004,
%   author  = {Gertler, Paul},
%   title   = {Do Conditional Cash Transfers Improve Child Health?},
%   journal = {American Economic Review},
%   year    = {2004},
%   volume  = {94},
%   number  = {2},
%   pages   = {336--341}
% }
%
% @article{parker2005,
%   author  = {Parker, Susan W. and Teruel, Graciela M.},
%   title   = {Randomization and Social Program Evaluation: The Case of {Progresa}},
%   journal = {The Annals of the American Academy of Political and Social Science},
%   year    = {2005},
%   volume  = {599},
%   pages   = {199--219}
% }
%
% @article{skoufias2001,
%   author  = {Skoufias, Emmanuel and Parker, Susan W. and Behrman, Jere R.
%              and Pessino, Carola},
%   title   = {Conditional Cash Transfers and Their Impact on Child Work and Schooling},
%   journal = {Econom\'ia},
%   year    = {2001},
%   volume  = {2},
%   number  = {1},
%   pages   = {45--96}
% }
%
% @article{barham2013,
%   author  = {Barham, Tania and Rowberry, Jacob},
%   title   = {Living longer: The effect of the {Mexican} conditional cash
%              transfer program on elderly mortality},
%   journal = {Journal of Development Economics},
%   year    = {2013},
%   volume  = {105},
%   pages   = {226--236},
%   doi     = {10.1016/j.jdeveco.2013.08.002}
% }
%
% @incollection{parker2008,
%   author    = {Parker, Susan W. and Rubalcava, Luis and Teruel, Graciela M.},
%   title     = {Evaluating conditional schooling-health transfer programs},
%   booktitle = {Handbook of Development Economics},
%   year      = {2008},
%   volume    = {4},
%   pages     = {3963--3035}
% }
%===========================================================================
